%% Description of the Paderborn bearing data dataset.
%% Covers dimensions, combinations of faults/bearings/conditions, EDA findings,
%% supported tests and splits; the test rig is not described here.
%% Build (from eda/paderborn/):   pdflatex paderborn && pdflatex paderborn

\documentclass[11pt,a4paper]{article}
\usepackage[T1]{fontenc}
\usepackage[utf8]{inputenc}
\usepackage[margin=2.4cm]{geometry}
\usepackage{amsmath}
\usepackage{booktabs}
\usepackage{graphicx}
\usepackage[hidelinks]{hyperref}

\title{Paderborn bearing data}
\author{}
\date{\today}

\begin{document}
\maketitle

\section{Source}

\section{Dimensions of the time series}

\section{Combinations of faults, bearings and operating conditions}

\begin{table}[h]
\centering
\small
\begin{tabular}{lp{5.2cm}p{5.2cm}}
\toprule
Class (Lessmeier) & Artificial damage & Real damage (accelerated lifetime test) \\
\midrule
Healthy    & \multicolumn{2}{l}{K001, K002, K003, K004, K005, K006} \\
Outer ring & KA01, KA03, KA05, KA06, KA07, KA08, KA09 & KA04, KA15, KA16, KA22, KA30, KB27$^\dagger$ \\
Inner ring & KI01, KI03, KI05, KI07, KI08 & KI04$^*$, KI14$^*$, KI16, KI17, KI18, KI21, KB23$^\dagger$, KB24$^\dagger$ \\
\bottomrule
\end{tabular}
\end{table}

\begin{itemize}
  \item Each entry is one physical bearing: 32 in total (6 healthy, 12 artificial, 14 real damage).
  \item Every bearing is recorded at all four operating conditions, 20 measurements each
        (128 runs, 2,560 measurements). One measurement: ca.\ 4\,s at 64\,kHz; vibration and
        two motor phase currents (D48).
  \item Operating conditions (speed, load torque, radial force):
        N15\_M07\_F10 (1500\,rpm, 0.7\,Nm, 1000\,N), N09\_M07\_F10 (900\,rpm, 0.7\,Nm, 1000\,N),
        N15\_M01\_F10 (1500\,rpm, 0.1\,Nm, 1000\,N), N15\_M07\_F04 (1500\,rpm, 0.7\,Nm, 400\,N).
  \item Artificial damage: EDM, drilling or electric engraver, extent level 1--2.
        Real damage: pitting or particle-caused indentations, extent level 1--3.
        Manufacturers IBU, MTK, FAG.
  \item Rows: the class Lessmeier et al.\ (2016) give each bearing (\texttt{lessmeier\_class}, D212).
  \item $^\dagger$ Combined damage on both rings (KB bearings), classed by the dominant damage
        (Lessmeier Sec.\ 7.4): KB23 pitting on both rings; KB24 distributed inner-ring pitting
        and outer-ring indentations; KB27 indentations on both rings.
  \item $^*$ Inner-ring pitting plus small outer-ring indentations (extent not given on the
        sheet); treated as plain inner-ring bearings (Lessmeier Sec.\ 7.3, D220).
  \item Healthy bearings differ by run-in time (1\,h to $>$50\,h).
  \item The authors' protocol (D220) uses 15 bearings, real damage only for the faults:
        K001--K005; KA04, KA15, KA16, KA22, KA30; KI04, KI14, KI16, KI18, KI21.
  \item Source: \texttt{BEARINGS}, \texttt{CONDITIONS} in \texttt{harness/data/paderborn.py}
        (from the damage sheets).
\end{itemize}

\section{Findings from EDA}

\section{Tests the data support}

\section{Splits}

\end{document}
